\setchapterpreamble[u]{\margintoc}
\chapter{Linear Algebra behind Learning from Data}
\labch{learning-from-data}

\sources{Strang, \emph{Introduction to Linear Algebra}, 6th ed.\ \cite{strang2023ila},
Chapter~10; Zhang, Lipton, Li and Smola, \emph{Dive into Deep Learning}
\cite{d2l2023}, Chapter~2.}

This chapter closes Part~I by putting the factorisations to work on data. Nothing
new is proved; the point is to see that the objects of the previous nine
chapters are the objects a learning system manipulates.

\section{Data as a matrix}
\labsec{data-matrix}

Store \(n\) examples with \(d\) features each as \(X\in\R^{n\times d}\), one
example per row. Two readings, both from \refsec{four-ways}, are used
constantly:

\begin{description}
	\item[By rows] each row is a point in \(\R^{d}\), and the dataset is a cloud
	of \(n\) such points.
	\item[By columns] each column is one feature measured across all examples,
	and the column space of \(X\) is the set of quantities expressible as
	weighted mixtures of features.
\end{description}

A linear model is then \(\hat{\vect{y}}=X\vect{w}\), so the achievable
predictions form exactly \(\Col(X)\), and fitting it by least squares is the
projection of \refch{orthogonality}.

\marginnote{This is the sense in which linear algebra is not a prerequisite for
learning but a description of it. The reachability question of
\refch{vectors-and-column-space} is the question of what a linear model can
represent.}

\section{Centring and covariance}
\labsec{covariance}

Subtracting the mean of each column gives the centred matrix \(X_0\). The
\emph{sample covariance matrix} is
\[
	S=\frac{1}{n-1}X_0\tp X_0\in\R^{d\times d},
\]
symmetric and positive semidefinite by \refsec{positive-definite}, with
\(s_{jj}\) the variance of feature \(j\) and \(s_{jk}\) the covariance between
features \(j\) and \(k\).

\begin{example}
Let
\[
	X_0=\begin{bmatrix}2&1\\1&2\\-2&-1\\-1&-2\end{bmatrix},
\]
whose columns already sum to zero. Then
\(X_0\tp X_0=\begin{bmatrix}10&8\\8&10\end{bmatrix}\) and
\(S=\tfrac13X_0\tp X_0\). The eigenvalues of \(X_0\tp X_0\) are \(18\) and \(2\),
with eigenvectors \([1,1]\tp\) and \([1,-1]\tp\).
\end{example}

\section{Principal components}
\labsec{pca}

Principal component analysis asks for the direction \(\vect{q}\) of unit length
along which the projected data vary most. That is the maximiser of
\(\vect{q}\tp S\vect{q}\), and by the spectral theorem it is the eigenvector of
\(S\) belonging to the largest eigenvalue.

\begin{proposition}
The principal directions of \(X_0\) are its right singular vectors, and the
variance captured by the \(i\)th component is proportional to \(\sigma_i^{2}\).
\end{proposition}

\begin{proof}
Writing \(X_0=U\Sigma V\tp\) gives \(X_0\tp X_0=V\Sigma\tp\Sigma V\tp\), which is
the eigendecomposition of a symmetric matrix with eigenvalues \(\sigma_i^{2}\)
and eigenvectors \(\vect{v}_i\).
\end{proof}

\begin{example}[continued]
Here \(\sigma_1=\sqrt{18}=3\sqrt2\) and \(\sigma_2=\sqrt2\), so the first
component \([1,1]\tp/\sqrt2\) accounts for \(18/20=90\%\) of the total variance.
Projecting onto it replaces two features by one while discarding a tenth of the
spread.
\end{example}

\marginnote{Computing the singular value decomposition of \(X_0\) is preferable
to forming \(S\) first, for the conditioning reason given in
\refsec{conditioning}.}

\section{Low rank as a modelling assumption}
\labsec{low-rank}

Eckart--Young (\refsec{eckart-young}) says the truncated sum
\(A_k=\sum_{i\le k}\sigma_i\vect{u}_i\vect{v}_i\tp\) is the best rank-\(k\)
approximation. Three familiar techniques are that theorem applied to different
matrices: dimensionality reduction applies it to a data matrix, compression of a
trained layer applies it to a weight matrix, and latent-factor recommendation
applies it to a partially observed ratings matrix.

\begin{remark}
The assumption being made each time is that the interesting structure lives in
few directions and the rest is noise. Whether that is true is a modelling
question about the data, not a mathematical one --- the theorem is silent about
where to cut.
\end{remark}

\section{Batched computation}
\labsec{batching}

A linear layer applied to a batch is one matrix product. With \(X\in\R^{n\times d}\)
and weights \(W\in\R^{d\times h}\),
\[
	Z=XW+\vect{1}\vect{b}\tp\in\R^{n\times h},
\]
where \(\vect{1}\in\R^{n}\) broadcasts the bias across rows. The block reading of
\refsec{four-ways} is what makes this efficient: the product decomposes into
independent blocks that hardware evaluates in parallel, which is the practical
reason mini-batches exist at all.

\begin{remark}
Bias makes the layer \emph{affine} rather than linear, so it violates
\(T(\vect{0})=\vect{0}\) from \refsec{linear-maps}. Nothing breaks --- an affine
map is a linear map followed by a shift --- but statements proved for linear
maps have to be applied to the linear part alone.
\end{remark}

\section{What carries forward}
\labsec{part1-summary}

Four ideas from Part~I are used without further comment in the rest of the book:
the column space as the set of reachable outputs; orthogonal projection as the
solution of a least-squares problem; the eigenvalues of a symmetric matrix as
the curvature of a quadratic; and the singular value decomposition as the
canonical way to separate signal from remainder. \refch{column-space-to-layers}
returns to them once the calculus and probability are in place.
